% TOP_PROBLEM: {"id": 3284, "title": "Complexity of the H-factor problem.", "category_id": 3, "category": {"id": 3, "name": "graph_theory", "display_name": "Graph Theory", "description": "Problems involving graphs, networks, and their properties.", "slug": "graph-theory", "order_index": 3, "created_at": "2026-07-31T15:26:25.671Z"}, "status": "open", "problem_number": "OPG-46738", "set_id": 12, "statusQualification": "Interprets the question as classification by fixed H and 0<c<1; a universal hardness reading is false, and c>=1 makes the simple-graph promise empty."}
% FIRST_POSTED: 2026-10-01
% ENTRY_KIND: statement_only
% UnsolvedMath ID 3284; source catalogue code OPG-46738
% !TeX program = lualatex
% Definition and sources only; this file is not an LLM solution attempt.
\documentclass[11pt,a4paper]{article}
\usepackage[margin=25mm]{geometry}
\usepackage{fontspec}
\setmainfont{lmroman10-regular.otf}[BoldFont=lmroman10-bold.otf,
 ItalicFont=lmroman10-italic.otf,BoldItalicFont=lmroman10-bolditalic.otf]
\usepackage{amsmath,amssymb,mathtools}
\usepackage{unicode-math}
\setmathfont{latinmodern-math.otf}
\AtBeginDocument{\let\setminus\smallsetminus}
\usepackage{xurl,hyperref}
\hypersetup{colorlinks=true,urlcolor=blue}
\setcounter{secnumdepth}{0}
\setlength{\parindent}{0pt}
\setlength{\parskip}{5pt}
\setlength{\emergencystretch}{3em}
\begin{document}
\section*{Complexity of the H-factor problem.}\label{3284}
{\small UnsolvedMath ID 3284 · OPG-46738 · Graph Theory · OpenGarden}
\subsection{Definitions and mathematical statement}
Fix a nonempty finite simple graph $H$ with $h$ vertices and a real density parameter $0<c<1$. An $H$-factor of a finite simple graph $G$ is a family of pairwise vertex-disjoint subgraphs isomorphic to $H$ whose vertex sets cover $V(G)$. Copies need not be induced: additional edges among or between their vertices may be ignored. Necessarily $h$ divides $|V(G)|$.

Consider the decision problem whose input is an $n$-vertex graph $G$ promised to satisfy $\delta(G)\ge cn$, and whose output specifies whether an $H$-factor exists. The graph $H$ and constant $c$ are fixed, rather than part of the growing input. The inequality means $\delta(G)\ge\lceil cn\rceil$. Inputs violating divisibility can be rejected immediately.

The source’s complexity question asks for which choices of $H$ and $c$ this promised problem is NP-hard, and for which choices a polynomial-time decision algorithm exists. Hardness is under polynomial-time many-one reductions whose output graphs satisfy the degree promise; running-time exponents and constants in an algorithm may depend on the fixed parameters.

This is a parameter-dependent classification, rather than a universal assertion of hardness for all $H,c$. For example, a one-edge factor is the perfect-matching problem, and sufficiently strong density hypotheses can guarantee factors in other cases. Values $c\ge1$ would give an empty simple-graph promise. The unresolved target includes the precise transition between algorithmic and hard regimes, with equality at a critical density treated separately when necessary.
\subsection{Short English statement}
Classify, for each fixed graph H and density c, the computational difficulty of finding whether a dense graph has a perfect packing by copies of H.
\subsection{Sources}
\begin{itemize}
\item[S1] ULAM.AI and the UnsolvedMath Contributors, supplied problems.json, record 3284 (OPG-46738), OpenGarden. Catalogue identity and source transcription; CC BY 4.0.
\url{https://huggingface.co/datasets/ulamai/UnsolvedMath}.
\item[S2] Open Problem Garden, Complexity of the H-factor problem.. Original problem and discussion.
\url{http://www.openproblemgarden.org/op/complexity_of_the_h_factor_problem}.
\item[S3] J. Han and A. Treglown, The complexity of perfect matchings and packings in dense hypergraphs.
\url{https://arxiv.org/abs/1609.06147}.
\end{itemize}
\end{document}
